\section{What this project has verified for itself}

\subsection{The combinatorial half, checked by the Lean kernel}

Our Lean development is written in core Lean~4, without Mathlib. It defines
primality, the twin prime conjecture, the predicate \texttt{BoundedGaps B}
(infinitely many pairs of distinct primes at distance at most $B$),
admissibility, and $\DHL[k,2]$. It contains no \texttt{sorry}, and the only
axioms it uses are \texttt{propext} and \texttt{Quot.sound}. The statement
$\DHL[k,2]$ is always a hypothesis, never an axiom.

\begin{theorem}[\lbl{Lean-checked}]\label{thm:lean}
\begin{enumerate}
\item The twin prime conjecture is equivalent to \texttt{BoundedGaps 2}.
\item If a finite Boolean check succeeds on a list $\mathcal H$ (it misses a
class modulo every $m$ with $2\le m\le|\mathcal H|$), then $\mathcal H$ is
admissible. The moduli $p>|\mathcal H|$ are handled by a formal pigeonhole
argument.
\item For each pair $(k,D)$ in Table~\ref{tab:tuples}, for
$(k,D)=(56,278)$, $(57,282)$, $(58,288)$, $(59,300)$, $(60,304)$, $(61,310)$,
$(62,320)$, $(63,324)$, $(64,330)$, and for $(102,578)$ and $(105,600)$, there is an explicit admissible $k$-tuple in $[0,D]$, and
$\DHL[k,2]$ implies \texttt{BoundedGaps D}. In particular $\DHL[2,2]$ implies
the twin prime conjecture.
\item $\DHL[k,2]$ implies $\DHL[k',2]$ for every $k'\ge k$
(\texttt{DHL.mono}: delete shifts from the tuple). So a proof of $\DHL[j,2]$
with $j\le k$ gives \texttt{BoundedGaps D} for each listed pair $(k,D)$.
\item \texttt{BoundedGaps 1} is false (\texttt{not\_boundedGaps\_one}), and a
pair $\{0,d\}$ is admissible if and only if $d$ is even
(\texttt{admissible\_pair\_iff}). So $2$ is the smallest bound of this form
that can hold, and the twin prime conjecture is the strongest statement of the
form \texttt{BoundedGaps D}.
\end{enumerate}
\end{theorem}

\begin{table}[h]
\caption{Pairs $(k,D)$ with a kernel-checked admissible $k$-tuple of diameter
$D$. The tuples were found by our own search; each $D$ agrees with the value
of $H(k)$ in published tables, whose optimality we have not re-proved.}
\label{tab:tuples}
\begin{tabular}{@{}rr@{\qquad}rr@{\qquad}rr@{\qquad}rr@{\qquad}rr@{\qquad}rr@{}}
\toprule
$k$&$D$&$k$&$D$&$k$&$D$&$k$&$D$&$k$&$D$&$k$&$D$\\
\midrule
2&2&11&36&20&80&29&130&38&176&47&226\\
3&6&12&42&21&84&30&136&39&182&48&236\\
4&8&13&48&22&90&31&140&40&186&49&240\\
5&12&14&50&23&94&32&146&41&188&50&246\\
6&16&15&56&24&100&33&152&42&196&51&252\\
7&20&16&60&25&110&34&156&43&200&52&254\\
8&26&17&66&26&114&35&158&44&210&53&264\\
9&30&18&70&27&120&36&162&45&212&54&270\\
10&32&19&76&28&126&37&168&46&216&55&272\\
\bottomrule
\end{tabular}
\end{table}

Theorem~\ref{thm:lean} is elementary. Its value is that implication
\eqref{eq:reduction} is now a checked interface: any future proof of
$\DHL[k,2]$, by anyone, yields the corresponding gap bound with nothing left
to trust on the combinatorial side.

\subsection{The sieve constant, with exact certificates}

Let $\R_k=\{t\in[0,1]^k:\ t_1+\dots+t_k\le1\}$ and, for square-integrable $F$
supported on $\R_k$,
\[
I_k(F)=\int_{\R_k}F^2,\qquad
J_k^{(m)}(F)=\int\left(\int F\,dt_m\right)^{2}\prod_{i\ne m}dt_i,\qquad
M_k=\sup_F\frac{\sum_{m=1}^{k}J_k^{(m)}(F)}{I_k(F)}.
\]
Maynard's theorem \cite[Prop.~4.2]{May15} states that if the primes have level
of distribution $\theta$ and $\theta M_k/2>1$ then $\DHL[k,2]$ holds. With
Bombieri--Vinogradov ($\theta<1/2$) the condition is $M_k>4$.

We take $F=\sum_{a+2b\le d}c_{a,b}(1-P_1)^aP_2^b$ with $P_j=\sum_it_i^j$, as
in \cite[\S8]{May15}. All integrals reduce to
\[
\int_{\R_k}(1-P_1)^AP_2^B\,dt=\frac{A!}{(k+A+2B)!}\,G_k(B),\qquad
\sum_{B\ge0}G_k(B)\frac{x^B}{B!}=\left(\sum_{b\ge0}\frac{(2b)!}{b!}x^b\right)^{k},
\]
which follows from the Dirichlet integral. The two quadratic forms are then
exact rational matrices. We compute an approximate top eigenvector in
floating point, round it to rationals, and evaluate the Rayleigh quotient in
exact arithmetic; the last step is a rigorous lower bound whatever happened
before it.

\begin{theorem}[\lbl{Computed}, exact rational arithmetic]\label{thm:Mk0}
With $F$ a polynomial in $1-P_1$ and $P_2$ only,
\[
M_{105}\ge4.00206976,\qquad M_{102}\ge4.00271343,\qquad M_{5}\ge2.00410541 .
\]
\end{theorem}

The first value reproduces Maynard's $M_{105}>4$ with $d=11$ ($42$
coefficients). The third gives $\DHL[5,2]$ and $H_1\le12$ under the
Elliott--Halberstam conjecture, again reproducing \cite{May15}. For
$d=11,15,19$ the least $k$ at which this family certifies $M_k>4$ is
$105,102,102$, and at $k=101$ it gives only $3.99931$ ($d=19$): polynomials in
$P_1,P_2$ alone stop at $k=102$.

\subsection{All symmetric test functions}
To go further we allow every symmetric polynomial. Since $P_1,\dots,P_k$
generate the symmetric polynomials, the functions
\[
(1-P_1)^a\,P^{b}:=(1-P_1)^a\prod_{j\ge2}P_j^{b_j},\qquad a+w(b)\le d,\quad w(b)=\sum_j jb_j,
\]
span all symmetric polynomials of degree at most $d$ (for $k\ge d$). The
integrals have a closed form that we have not seen stated in this way
(\lbl{Proved here}; possibly folklore).

\begin{lemma}\label{lem:egf}
Let $X_1,\dots,X_k$ be independent exponential random variables of mean $1$
and $E_k(b)=\mathbb E\prod_jP_j(X)^{b_j}$. Then
\[
\int_{\R_k}(1-P_1)^AP^{b}\,dt=\frac{A!}{(k+A+w(b))!}\,E_k(b),
\qquad
\sum_b E_k(b)\frac{s^b}{b!}=\left(\sum_{c}w(c)!\,\frac{s^c}{c!}\right)^{k},
\]
where $s^b=\prod_js_j^{b_j}$ and $b!=\prod_jb_j!$. In particular each
$E_k(b)$ is a positive integer, computed exactly by multiplying
exponential generating functions.
\end{lemma}

\begin{proof}
For a monomial, the Dirichlet integral gives
$\int_{\R_k}(1-P_1)^At^\alpha\,dt=A!\,\alpha!/(k+A+|\alpha|)!$, and
$\alpha!=\mathbb EX^\alpha$. Summing over the monomials of $P^b$, all of
degree $w(b)$, gives the first identity. For the second,
$\mathbb E\exp(\sum_js_jP_j(X))=\left(\mathbb E\exp(\sum_js_jX_1^j)\right)^k$
by independence, and
$\mathbb E\exp(\sum_js_jX_1^j)=\sum_c\mathbb E X_1^{w(c)}\,s^c/c!=\sum_cw(c)!\,s^c/c!$,
as formal power series.
\end{proof}

The $J$-functional reduces to the same lemma in $k-1$ variables: with
$P_j=P_j'+t_k^j$,
\[
\int_0^{1-P_1'}(1-P_1)^aP^b\,dt_k=\sum_{e\le b}\binom{b}{e}
\frac{a!\,w(e)!}{(a+w(e)+1)!}\,(1-P_1')^{a+w(e)+1}P'^{\,b-e}.
\]
So both Gram matrices are exact rational matrices. We find an approximate top
eigenvector by inverse iteration with a fixed shift above
$\frac{k}{k-1}\log k$, which can only converge to the top eigenvalue. We take
its floating-point midpoints as exact dyadic rationals and evaluate the
Rayleigh quotient exactly. The working precision of the eigenvector matters
only for speed. A separate script, \texttt{code/verify\_cert.py}, rebuilds
both Gram matrices from scratch and re-evaluates each stored coefficient
vector in exact arithmetic.

\begin{theorem}[\lbl{Computed}, exact rational arithmetic]\label{thm:Mk}
\begin{gather*}
M_{105}\ge4.17427,\qquad M_{64}\ge4.00363,\qquad M_{58}\ge4.00927,\\
M_{56}\ge4.00519,\qquad M_{55}\ge4.00516,\qquad M_{54}\ge4.00238 .
\end{gather*}
The first two use all symmetric polynomials of degree $11$ and $12$
($195$ and $272$ coefficients). The others use the subspace spanned by
$(1-P_1)^aP^b$ with $b$ supported on even $j$, as in \cite[\S7]{Pol14b}, of
degree $16$, $18$, $20$ and $23$ ($307$, $471$, $707$ and $1236$ coefficients).
\end{theorem}

\begin{center}
\begin{tabular}{@{}lccc@{}}
\toprule
Basis & degree & coefficients & least $k$ with certified $M_k>4$\\
\midrule
$1-P_1$, $P_2$ & 19 & 110 & 102\\
all symmetric & 12 & 272 & 64\\
even power sums & 16 & 307 & 58\\
even power sums & 18 & 471 & 56\\
even power sums & 20 & 707 & 55$^{\dagger}$\\
even power sums & 23 & 1236 & 54$^{\dagger}$\\
\bottomrule
\end{tabular}
\end{center}
\noindent{\small $^{\dagger}$ Only this value of $k$ was run at this degree; smaller $k$ were not tried.}

\begin{corollary}[end-to-end within this project]\label{cor:578}
Assuming the Bombieri--Vinogradov theorem and \cite[Prop.~4.2]{May15}, both
\lbl{Literature}, Theorems~\ref{thm:Mk} and~\ref{thm:lean} give
$H_1\le270$.
\end{corollary}

This reproduces a known result and is not new. Polymath8b obtained
$M_{54}>4.00238$ with even-part polynomials of degree $23$
\cite[Thm.~3.9]{Pol14b}, and hence $H_1\le270$. Our certified value in the
same space, $4.002386\ldots$, agrees with theirs to every digit they print,
which is an external check on our pipeline; two independent runs of it (one
from an earlier session) produced the same value. Corollary~\ref{cor:578}
is the best that the plain simplex gives with Bombieri--Vinogradov (see the
next subsection); Corollary~\ref{cor:246} below takes the next rung of the
ladder $600\to270\to246\to240\to212\to186\to182$.

\subsection{What the plain method cannot do}
Polymath8b proved $M_k\le\frac{k}{k-1}\log k$ \cite{Pol14b}. Hence
(\lbl{Computed} from that bound) $M_k>4$ forces $k\ge51$, and on the plain
simplex $\DHL[k,2]$ requires a level of distribution
$\theta>2(k-1)/(k\log k)$:
\begin{center}
\begin{tabular}{@{}rcccccccc@{}}
\toprule
$k$ & 50 & 45 & 40 & 35 & 30 & 20 & 10 & 5\\
\midrule
$\theta>$ & 0.501 & 0.514 & 0.529 & 0.547 & 0.568 & 0.634 & 0.782 & 0.994\\
\bottomrule
\end{tabular}
\end{center}
The 2026 results go below $k=50$ by changing the support of $F$, not by
proving a uniform level of distribution of this size.

\subsection{Beyond the plain simplex: \texorpdfstring{$H_1\le246$}{H1 <= 246}}
Polymath8b's $\varepsilon$-enlarged simplex \cite[Thm.~3.12]{Pol14b}
(\lbl{Literature}; we re-read the statement in the arXiv text for this
version) takes $F$ supported on $(1+\varepsilon)\R_k$ and integrates the
square of each marginal $\int F\,dt_m$ only over $(1-\varepsilon)\R_{k-1}$.
Call the resulting constant $M_{k,\varepsilon}$. If $\mathrm{EH}[\theta]$
holds with $1+\varepsilon<1/\theta$ and $M_{k,\varepsilon}>2/\theta$, then
$\DHL[k,2]$ holds. Bombieri--Vinogradov gives every $\theta<1/2$, so for
every fixed $\varepsilon<1$ the condition is $M_{k,\varepsilon}>4$.

We extended our pipeline to this functional, following the route that
\cite[\S7.2]{Pol14b} describes in words. Put $L=1+\varepsilon$,
$\tau=1-\varepsilon$ and use the basis $(L-P_1)^aP^b$ ($b$ supported on even
$j$). The following identities are elementary (\lbl{Proved here}): the
substitution $t=Ls$ and Lemma~\ref{lem:egf} give
\[
\int_{L\R_k}(L-P_1)^AP^b\,dt=L^{k+A+w(b)}\,\frac{A!}{(k+A+w(b))!}\,E_k(b);
\]
the inner integral over $t_m\in[0,L-P_1']$ is the formula displayed after
Lemma~\ref{lem:egf} with $1-P_1'$ replaced by $L-P_1'$; and on
$\tau\R_{k-1}$ we write $L-P_1'=2\varepsilon+(\tau-P_1')$, so that
\[
\int_{\tau\R_{k-1}}(L-P_1')^AP'^{\,B}=\sum_{i=0}^{A}\binom Ai(2\varepsilon)^{A-i}\,
\tau^{k-1+i+w(B)}\,\frac{i!}{(k-1+i+w(B))!}\,E_{k-1}(B).
\]
For rational $\varepsilon$ both Gram matrices are again exact rational
matrices. We checked the code against hand computation for $k=2,3$ and
$F\in\{1,L-P_1\}$, and against direct iterated symbolic integration (sympy,
without Lemma~\ref{lem:egf} or the binomial rewrite) of every Gram entry for
$(k,d,\varepsilon)=(3,3,1/5)$ and $(4,4,1/25)$: no mismatch
(\texttt{code/check\_eps\_sympy.py}). At $\varepsilon=0$ it is the old code
path (an old certificate re-verifies unchanged). We also found that the working precision of
the eigenvector search can be cut from about $6000$ to $300$--$400$ bits
without changing any certified digit, which makes degree $25$ feasible on a
four-core machine (build $16$ minutes, search $4$ minutes).

\begin{theorem}[\lbl{Computed}, exact rational arithmetic]\label{thm:Meps}
With even power sums of degree $25$ ($1780$ coefficients) and of degree $27$
($2526$ coefficients),
\[
M_{50,1/25}\ge4.0012466\qquad\text{and}\qquad M_{50,1/25}\ge4.0043063 .
\]
\end{theorem}

The certificates are \texttt{code/cert\_M50e25\_D25e.json} and
\texttt{code/cert\_M50e25\_D27e.json} (the degree-$27$ run took $43$ minutes
to build and $10$ minutes to search, with about $10$~GB of memory); the script
\texttt{code/verify\_cert.py} rebuilds both Gram matrices from scratch and
re-evaluates the stored rational vector exactly; it uses the same
Gram-matrix builder, so it checks the arithmetic and the bookkeeping, while
the formulas themselves are checked by the symbolic test above. Polymath8b report
$M_{50,1/25}>4.00124$ at degree $25$ and $>4.0043$ at degree $27$
\cite[\S7.2]{Pol14b}; our values agree with both to every printed digit.


\begin{corollary}[end-to-end within this project]\label{cor:246}
Assuming the Bombieri--Vinogradov theorem and \cite[Thm.~3.12(i)]{Pol14b},
both \lbl{Literature}, Theorem~\ref{thm:Meps} and the Lean tuple $(50,246)$
of Theorem~\ref{thm:lean} give $H_1\le246$.
\end{corollary}

This is the established bound of 2014, reproduced; it is not new. It closes
the second rung of our ladder. Table~\ref{tab:eps} shows how the certified
values grow with the degree and depend on $\varepsilon$; it is the data
behind Remark~\ref{rem:k49} below.

\begin{table}[h]
\caption{Certified lower bounds for $M_{k,\varepsilon}$ (\lbl{Computed}, exact
Rayleigh quotients, even power sums of degree $D$). Data:
\texttt{runs/scan\_k49.jsonl}, \texttt{runs/scan\_k50.jsonl}, written by
\texttt{code/scan\_eps.py}; the column $k=50$, $D=16$ from direct runs of
\texttt{code/mk\_general.py}. A dash: not computed.}
\label{tab:eps}
\begin{tabular}{@{}lllllll@{}}
\toprule
& \multicolumn{3}{c}{$k=49$} & \multicolumn{3}{c}{$k=50$}\\
$\varepsilon$ & $D=16$ & $D=18$ & $D=20$ & $D=16$ & $D=18$ & $D=20$\\
\midrule
$0$    & 3.88640 & 3.90102 & 3.91056 & 3.90173 & 3.91718 & 3.92739\\
$1/50$ & 3.93491 & 3.95066 & 3.96110 & 3.94979 & 3.96644 & 3.97761\\
$1/30$ & 3.93969 & 3.95745 & 3.96933 & --      & 3.97254 & 3.98522\\
$1/25$ & 3.93674 & 3.95563 & 3.96840 & 3.95041 & 3.97034 & 3.98395\\
$1/20$ & 3.92741 & 3.94823 & 3.96244 & --      & 3.96234 & 3.97745\\
\bottomrule
\end{tabular}
\end{table}

